\chapter{Pruebas y validación}
\label{chap:pruebas}

\ph{Este capítulo describe cómo se ha verificado que el sistema cumple los
requisitos. Las secciones siguen el orden clásico de la \emph{pirámide de
pruebas}: de la base (muchas pruebas, rápidas y baratas) a la cumbre (pocas,
lentas y caras). En la base están las pruebas unitarias, en el medio las
de integración y en la cima las de extremo a extremo / aceptación.
Delimita el alcance desde el principio: este capítulo responde a la pregunta
de la \emph{verificación} («¿hemos construido bien el sistema?»); la
\emph{evaluación} de los resultados frente a los objetivos del trabajo y su
discusión se realiza en el Capítulo~\ref{chap:resultados}.}

\section{Estrategia de pruebas}
\label{sec:estrategia_pruebas}

\ph{Resumen del enfoque adoptado: niveles de pruebas usados, framework
elegido y entorno en el que se ejecutan. Adapta los niveles al tipo de TFG
(software, hardware, modelo de IA, estudio empírico\dots).}

\ph{Niveles de la pirámide aplicados en este TFG (incluye los que apliquen):}

\begin{itemize}
  \item \textbf{\ph{Pruebas unitarias}}: \ph{base de la pirámide. Verifican
        unidades aisladas (funciones, clases, módulos). Son las más numerosas
        y rápidas; deben ser deterministas y no depender de servicios externos.}
  \item \textbf{\ph{Pruebas de integración}}: \ph{nivel intermedio. Comprueban
        que varios componentes funcionan juntos correctamente, incluyendo
        contratos con servicios externos (BD, APIs internas\dots).}
  \item \textbf{\ph{Pruebas de extremo a extremo / sistema}}: \ph{cumbre de la
        pirámide. Validan flujos completos desde la perspectiva del usuario o
        cliente. Son las menos numerosas porque son las más lentas y frágiles.}
  \item \textbf{\ph{Pruebas de aceptación}} (opcional): \ph{cotejan el sistema
        frente a los requisitos del Capítulo~\ref{chap:requisitos}.}
\end{itemize}

\ph{Comenta brevemente el entorno de pruebas: si las dependencias externas
son reales, simuladas o efímeras (en memoria, contenedores desechables,
mocks\dots) y cómo se garantiza el aislamiento entre tests.}

\section{Pruebas unitarias}
\label{sec:pruebas_unitarias}

\ph{Describe cómo se prueban los componentes centrales del sistema en
aislamiento (la lógica de negocio, el algoritmo principal o el módulo más
crítico del TFG). Indica qué grupos de casos cubre, no la lista exhaustiva.}

\ph{Grupos de casos típicos:}

\begin{itemize}
  \item \textbf{\ph{Caso nominal}}: \ph{entrada válida que debe producir el
        resultado esperado.}
  \item \textbf{\ph{Entradas inválidas}}: \ph{datos mal formados, fuera de
        rango o con campos ausentes.}
  \item \textbf{\ph{Casos límite}}: \ph{valores en la frontera del dominio
        admisible.}
  \item \textbf{\ph{Reglas de negocio}}: \ph{combinaciones que activan
        restricciones específicas del dominio.}
\end{itemize}

\ph{Si la lógica del módulo no se reduce a un booleano, comenta también qué
información complementaria se valida (mensajes, códigos de error, listas de
campos afectados\dots).}

\section{Pruebas de integración}
\label{sec:pruebas_integracion}

\ph{Describe las pruebas que ejercitan varios componentes a la vez.
Conviene separarlas según los componentes que combinan.}

\begin{itemize}
  \item \ph{Pruebas sobre el contrato con la base de datos / capa de
        persistencia (ORM, esquemas, transacciones).}
  \item \ph{Pruebas sobre la capa de servicios o API
        (códigos de respuesta, contratos de entrada/salida, autenticación,
        control de acceso entre usuarios).}
  \item \ph{Pruebas sobre integraciones con servicios externos
        (proveedores de identidad, almacenamiento, APIs de terceros\dots).}
\end{itemize}

\ph{Documenta también las pruebas de casos negativos
(valores fuera de rango, restricciones violadas, datos incompletos\dots).}

\section{Pruebas de extremo a extremo}
\label{sec:pruebas_e2e}

\ph{Describe la validación del sistema desde la perspectiva del usuario
final. No tiene por qué ejecutarse en un navegador o entorno gráfico real:
muchas veces basta con automatizar el flujo a través de la interfaz pública
del sistema (HTTP, CLI, etc.). Estas pruebas son las menos numerosas pero
las que dan más confianza en que el sistema funciona como un todo.}

\ph{Plantilla de pasos del flujo principal validado:}

\begin{enumerate}
  \item \ph{Paso 1.}
  \item \ph{Paso 2.}
  \item \ph{Paso 3.}
  \item \ph{Paso 4.}
  \item \ph{Paso 5.}
\end{enumerate}

\ph{Casos complementarios habituales:}

\begin{itemize}
  \item \ph{Estado inicial vacío.}
  \item \ph{Repetición del flujo con varias instancias y comprobación del orden.}
  \item \ph{Acceso no autorizado o sesión inválida.}
  \item \ph{Recuperación tras un error controlado.}
\end{itemize}

\section{Cobertura y métricas de calidad}
\label{sec:cobertura}

\ph{Indica las métricas que has usado para evaluar la calidad de la suite
y, si procede, los umbrales mínimos exigidos en CI.}

\ph{Métricas habituales:}

\begin{itemize}
  \item \ph{Cobertura de líneas / ramas.}
  \item \ph{Número de pruebas y desglose por nivel de la pirámide.}
  \item \ph{Resultado del análisis estático (linter, tipado).}
\end{itemize}

\ph{Comenta también qué se excluye de la medición
(ficheros generados, plantillas, código de migración, dependencias externas\dots).}

\ph{Resume la suite completa en la Tabla~\ref{tab:resumen_suite}: el tribunal
debe poder ver de un vistazo cuántas pruebas hay, de qué tipo y con qué
resultado.}

\begin{table}[H]
\centering
\small
\begin{tabular}{@{}P{4.5cm} P{3cm} P{3cm} P{3cm}@{}}
\toprule
\textbf{Nivel} & \textbf{N.º de pruebas} & \textbf{Resultado} & \textbf{Cobertura} \\
\midrule
Unitarias         & \ph{N} & \ph{N/N pasan} & \ph{X\,\%} \\
Integración       & \ph{N} & \ph{N/N pasan} & \ph{X\,\%} \\
Extremo a extremo & \ph{N} & \ph{N/N pasan} & \ph{---} \\
\midrule
\textbf{Total}    & \textbf{\ph{N}} & \textbf{\ph{N/N pasan}} & \textbf{\ph{X\,\%}} \\
\bottomrule
\end{tabular}
\caption{Resumen de la suite de pruebas por nivel de la pirámide.}
\label{tab:resumen_suite}
\end{table}

\ph{No te limites a enumerar cifras: interprétalas. Una cobertura del X\,\%
no dice nada por sí sola. Indica qué módulos quedan menos cubiertos y por qué
es asumible (código trivial, generado, de difícil automatización\dots), y qué
te dice la distribución de pruebas por nivel sobre la fiabilidad de la suite
(¿la base de la pirámide es realmente la más ancha?).}

\ph{Cierra el eslabón con los objetivos: indica qué criterios de la
Tabla~\ref{tab:criterios-exito} quedan verificados por esta suite,
apoyándote en la matriz de la Tabla~\ref{tab:trazabilidad_rf_pruebas}.}

% --- Transición al siguiente capítulo ---
\ph{Cierra el capítulo con un párrafo puente de 2--4 líneas: qué queda
establecido aquí y qué pregunta abre el capítulo siguiente. Por ejemplo:
«Verificado que el sistema está bien construido y que cumple los requisitos
especificados, queda la pregunta de fondo: ¿resuelve el problema planteado y
en qué medida? El Capítulo~\ref{chap:resultados} evalúa los resultados
obtenidos frente a los criterios de éxito y los discute críticamente».}
